% !TeX TS-program = pdflatex

\documentclass[a4paper]{article}

% \usepackage[default]{fontsetup}

\usepackage{fancyhdr}
\usepackage{extramarks}
\usepackage{amsmath}
\usepackage{amsthm}
\usepackage{amsfonts}
\usepackage{tikz}
\usepackage[plain]{algorithm}
\usepackage{algpseudocode}
\usepackage{enumerate}
\usepackage{tikz}

\usetikzlibrary{automata,positioning}

%
% Basic Document Settings
%  

\topmargin=-0.2in
\evensidemargin=0in
\oddsidemargin=0in
\textwidth=6.5in
\textheight=9.5in
\headsep=0.25in

\linespread{1.1}

\pagestyle{fancy}
\lhead{\hmwkAuthorName}
\chead{\hmwkClass : \hmwkTitle}
\rhead{\firstxmark}
\lfoot{\lastxmark}
\cfoot{\thepage}

\renewcommand\headrulewidth{0.4pt}
\renewcommand\footrulewidth{0.4pt}

\setlength\parindent{0pt}

%
% Create Problem Sections
%

%修改下标的地方
%\newcommand{\enterProblemHeader}[1]{
	%	\nobreak\extramarks{}{Problem \arabic{#1} continued on next page\ldots}\nobreak{}
	%	\nobreak\extramarks{Problem \arabic{#1} (continued)}{Problem \arabic{#1} continued on next page\ldots}\nobreak{}
	%}

\newcommand{\exitProblemHeader}[1]{
    \nobreak\extramarks{Problem \arabic{#1} }{Problem \arabic{#1} continued on next page\ldots}\nobreak{}
	\stepcounter{#1}
	\nobreak\extramarks{Problem \arabic{#1}}{}\nobreak{}
}

\newcommand*\circled[1]{\tikz[baseline=(char.base)]{
		\node[shape=circle,draw,inner sep=2pt] (char) {#1};}}


\setcounter{secnumdepth}{0}
\newcounter{partCounter}
\newcounter{homeworkProblemCounter}
\setcounter{homeworkProblemCounter}{1}
\nobreak\extramarks{Problem \arabic{homeworkProblemCounter}}{}\nobreak{}

%
% Homework Problem Environment
%
% This environment takes an optional argument. When given, it will adjust the
% problem counter. This is useful for when the problems given for your
% assignment aren't sequential. See the last 3 problems of this template for an
% example.
%

\newenvironment{homeworkProblem}[1][-1]{
	\ifnum#1>0
	\setcounter{homeworkProblemCounter}{#1}
	\fi
	\section{Problem \arabic{homeworkProblemCounter} }
	\setcounter{partCounter}{1}
	% \enterProblemHeader{homeworkProblemCounter} % 注释掉未定义的命令
}{
	\exitProblemHeader{homeworkProblemCounter}
}

%
% Homework Details
%   - Title
%   - Class
%   - Due date
%   - Name
%   - Student ID

\newcommand{\hmwkTitle}{Homework\ \#05}
\newcommand{\hmwkClass}{Probability \& Statistics for EECS}
\newcommand{\hmwkDueDate}{2026-11-3}
\newcommand{\hmwkAuthorName}{}
\newcommand{\hmwkAuthorID}{}


%
% Title Page
%

\title{
	\vspace{2in}
	\textmd{\textbf{\hmwkClass:\\  \hmwkTitle}}\\
	\normalsize\vspace{0.1in}\small{Due\ on\ \hmwkDueDate\ at 23:59}\\
	\vspace{4in}
}

\author{
	Name: \textbf{\hmwkAuthorName} \\
	Student ID: \hmwkAuthorID}
\date{}

\renewcommand{\part}[1]{\textbf{\large Part \Alph{partCounter}}\stepcounter{partCounter}\\}

%
% Various Helper Commands
%

% Useful for algorithms
\newcommand{\alg}[1]{\textsc{\bfseries \footnotesize #1}}
% For derivatives
\newcommand{\deriv}[1]{\frac{\mathrm{d}}{\mathrm{d}x} (#1)}
% For partial derivatives
\newcommand{\pderiv}[2]{\frac{\partial}{\partial #1} (#2)}
% Integral dx
\newcommand{\dx}{\mathrm{d}x}
% Alias for the Solution section header
\newcommand{\solution}{\textbf{\large Solution}}
% Probability commands: Expectation, Variance, Covariance, Bias
\newcommand{\E}{\mathrm{E}}
\newcommand{\Var}{\mathrm{Var}}
\newcommand{\Cov}{\mathrm{Cov}}
\newcommand{\Bias}{\mathrm{Bias}}

\begin{document}
	
	
	% \maketitle
	% \thispagestyle{empty}
	% \pagebreak
	
	\date{
		Due on Nov. 3, 2026, 18:00 UTC+8}
	\title{SI 140A-01  Probability \& Statistics for EECS, Fall 2026 \\
		Homework 5}
	\maketitle
	Read all the instructions below carefully before you start working on the assignment, and before you make a submission.
	\begin{itemize}
		\item You are required to write down all the major steps towards making your conclusions; otherwise you may obtain limited points of the problem.
		\item Write your homework in English; otherwise you will get no points of this homework.
		\item Any form of plagiarism will lead to $0$ point of this homework. 
	\end{itemize}
	
	
	
	\newpage
  
	\begin{homeworkProblem}
	 If $X \sim \text{HGeom}(w,b,n)$, what is the distribution of $n-X$? Give a short proof.



   
		
	\end{homeworkProblem}
	
	\newpage
	\begin{homeworkProblem}
		     A particular Youtuber is known for his arbitrary steak-eating habit during live streaming.
    In each live streaming day, he orders a steak with doneness as one of well done, medium well, medium, medium rare, and rare (with equal probability, independent of other live streaming days).
    How many days of live streaming do you expect to watch before you see him eating steaks with each possible doneness at least one once (suppose you are a big fan who watches every live streaming)?

	\end{homeworkProblem}
	
	
	\newpage
	\begin{homeworkProblem}
	  Mario and Zelda are independently performing independent Bernoulli trials.
    For concreteness, assume that Mario is flipping a nickel with probability $p_{1}$ of Heads and Zelda is flipping a penny with probability $p_2$ of Heads.
    Let $X_{1}, X_{2}, \cdots $ be Mario's results and $Y_{1}, Y_{2}, \cdots $ be Zelda's results, with $X_{i} \sim \text{Bern}(p_{1})$ and $Y_{j} \sim \text{Bern}(p_{2})$.
    
    \begin{enumerate}[(a)]
        \item
        Find the distribution and expected value of the first time at which they are simultaneously successful, \textit{i.e.}, the smallest $n$ such that $X_{n} = Y_{n} = 1$.
        
        \item
        Find the expected time until at least one has a success (including the success).
        
        \item
        For $p_{1} = p_{2}$, find the probability that their first successes are simultaneous, and use this to find the probability that Mario's first success precedes Zelda's.
    \end{enumerate}

\end{homeworkProblem}


\newpage
\begin{homeworkProblem} 
Show that for any events $A_1,\dots,A_n$,
\[
P(A_1 \cap A_2 \cap \cdots \cap A_n) \ge \sum_{j=1}^n P(A_j) - n + 1.
\]

Hint: First prove a similar statement about indicator r.v.s, by interpreting what the events $I(A_1 \cap A_2 \cdots \cap A_n)=1$ and $I(A_1 \cap A_2 \cdots \cap A_n)=0$ mean.

\end{homeworkProblem}

	
\end{document}